\chapter{Dasadhamma Sutta}

Bhikkhus, há dez dhammas sobre os quais aquele que se ordenou deve reflectir repetidamente. Quais são esses dez?

\begin{enumerate}
\item
  Já não vivo de acordo com objectivos e valores mundanos. Isto deve ser meditado sabiamente por quem se dedicou à vida monástica.
\item
  A minha própria vida é sustentada pelas doações dos outros. Isto deve ser meditado sabiamente por aquele que se ordenou.
\item
  Devo esforçar"-me por abandonar os meus hábitos anteriores. Isto deve ser meditado sabiamente por aquele que se ordenou.
\item
  Surge na minha mente algum sabiamente pela minha conduta? Isto deve ser reflectido sabiamente por quem se dedicou à vida monástica.
\item
  Será que os meus companheiros espirituais poderiam criticar a minha conduta? Isto deve ser reflectido sabiamente por quem se dedicou à vida monástica.
\item
  Tudo o que é meu, amado e agradável, tornar"-se"-á diferente, separar"-se"-á de mim. Isto deve ser reflectido sabiamente por aquele que se ordenou.
\item
  Eu sou o dono do meu kamma, herdeiro do meu kamma, nascido do meu kamma, relacionado com o meu kamma, permaneço sustentado pelo meu kamma; qualquer kamma que eu venha a fazer, para o bem ou para o mal, disso serei o herdeiro. Isto deve ser meditado sabiamente por aquele que se dedicou à vida monástica.
\item
  Os dias e as noites passam implacavelmente, como estou a passar o meu tempo? Isto deve ser meditado sabiamente por aquele que se ordenou.
\item
  Gosto da solidão ou não? Isto deve ser reflectido sabiamente por aquele que se ordenou.
\item
  A minha prática deu frutos em forma de liberdade ou discernimento, de modo que, no fim da minha vida, eu não precise sentir vergonha quando questionado pelos meus? Isto deve ser reflectido sabiamente por aquele que se ordenou.
\end{enumerate}

Bhikkhus, estes são os dez dhammas sobre os quais aquele que se ordenou deve reflectir sabiamente.

\emph{Aṅguttara Nikāya 10.48}
